% !TITLE 雨霖铃·寒蝉凄切
% !AUTHOR 柳永
% !DYNASTY 两宋
% !GENRE 词
% !ORDER 4

\begin{poem}
寒蝉\textsuperscript{1}凄切，对长亭\textsuperscript{2}晚，骤雨初歇。
都门帐饮\textsuperscript{3}无绪\textsuperscript{4}，留恋处，兰舟\textsuperscript{5}催发。
执手相看泪眼，竟无语凝噎\textsuperscript{6}。
念去去\textsuperscript{7}，千里烟波，暮霭\textsuperscript{8}沉沉楚天\textsuperscript{9}阔。

多情自古伤离别，更那堪、冷落清秋节\textsuperscript{10}！
今宵酒醒何处？杨柳岸、晓风残月\textsuperscript{11}。
此去经年\textsuperscript{12}，应是良辰好景虚设。
便纵有千种风情\textsuperscript{13}，更与何人说？
\end{poem}

\begin{annotations}
\notetext{1}{寒蝉：秋后的蝉。蝉的生命将尽，叫声凄凉，故称“寒蝉”。}
\notetext{2}{长亭：古代设在路边的亭舍，用作送别的场所，常是“十里一长亭，五里一短亭”。}
\notetext{3}{都门帐饮：在京城门外搭帐篷饯行饮酒。}
\notetext{4}{无绪：没有心情。虽设帐饯行，却心烦意乱，无心吃喝。}
\notetext{5}{兰舟：木兰木做的小船，是船的美称。也有人说指划船的船夫（鲁班曾刻木兰为舟）。}
\notetext{6}{凝噎：喉咙被堵住说不出话。噎（yē），气逆哽塞。太难过以至于一个字都说不出来。}
\notetext{7}{去去：走了又走，越走越远。叠用“去”字，写出路的漫长和别离的深切。}
\notetext{8}{暮霭：傍晚的雾气。霭（ǎi），云雾。}
\notetext{9}{楚天：楚地的天空。柳永从汴京（河南）出发南下，南方的天空开阔无边，衬托人的渺小和孤单。}
\notetext{10}{清秋节：清冷的秋季时节。离别本已令人难过，再逢冷落凄清的秋天，更加不堪。}
\notetext{11}{杨柳岸、晓风残月：杨柳岸边，晨风吹拂，天上挂着一弯残月——这是全词最有名的句子。杨柳代表离别（古人有折柳送别的习俗），残月代表残缺与孤寂。}
\notetext{12}{经年：年复一年，一年又一年。此去一别，不知多少年才能相见。}
\notetext{13}{风情：情意、柔情蜜意。即使有千万种深情，又能对谁说呢？}
\end{annotations}

\begin{appreciation}
《雨霖铃》是中国文学史上写离别最有名的一首词之一。“多情自古伤离别”——这句被后人引用上千年的词，就是从这首词里来的。

“寒蝉凄切，对长亭晚，骤雨初歇。”秋天的蝉、傍晚的长亭、刚停的雨，三个镜头，全是冷色。送行的酒席摆在京城门外，“都门帐饮无绪”——谁也没有心思吃喝。最难受的是“留恋处，兰舟催发”：舍不得走，船夫催了。时间不等人，再舍不得也得开。

于是有了全词最动人的画面：“执手相看泪眼，竟无语凝噎。”不是没有话说，是话太多，全堵在喉咙口。

“念去去，千里烟波，暮霭沉沉楚天阔。”船开了，暮霭沉沉，天地开阔，人小得像一粒沙。下阕推开一步：“多情自古伤离别，更那堪、冷落清秋节。”古往今来，多情的人都怕离别，何况还是冷清的秋天。

“今宵酒醒何处？杨柳岸、晓风残月。”酒醒时人在哪里？没有答案，只有杨柳岸、晨风、一弯残月。杨柳是折柳送别的柳——“柳”与“留”同音，折柳就是留人；残月是残缺的月。九个字里没有一个“愁”字，却比什么愁都凉。

“此去经年，应是良辰好景虚设。便纵有千种风情，更与何人说？”往后的好日子没有你，都是摆设；千种风情，又跟谁说。

还记得《行行重行行》吗？那首诗里我预告过，柳永写“衣带渐宽终不悔”，接的就是“衣带日已缓”——写这两句的柳永，正是这首《雨霖铃》的作者。《行行重行行》写等人的苦，《雨霖铃》写走的人的苦，合起来才是离别完整的样子。还有《少司命》里的“悲莫悲兮生别离”——这首词就是给这句古话做了一整套插图。

人生躲不开离别：升学、工作、搬家，每一次都是“兰舟催发”。但好好说再见，然后各自好好吃饭，该见面的时候自然会见。古人要杨柳和残月陪着过这个坎，我们不必。
\end{appreciation}
